\documentclass[a4paper, 12pt]{article}

\usepackage[utf8]{inputenc}
\usepackage[T1]{fontenc}
\usepackage[portuguese]{babel}
\usepackage{graphicx}
\usepackage{float}
% \usepackage{hyperref}
\usepackage[hidelinks]{hyperref}

\begin{document}

\renewcommand{\arraystretch}{1.25}

\noindent
\begin{tabular}{|p{\dimexpr\linewidth-2\tabcolsep-1.2pt}|}
    \hline
    \large WebAcademy\\
    \hline
    \large \textbf{Disciplina: }Computação em Nuvem - Resumo\\
    \hline
    \large \textbf{Nome: }Marcos Vinicius Souza da Costa\\
    \hline
\end{tabular}
 
\section{Virtualização: conceitos e exemplos. Clusters e data centers. }
A virtualização é um processo que permite utilizar os recursos de uma máquina física, como um servidor ou computador pessoal, de forma mais otimizada através abstração e particionamento dos recursos de hardware em máquinas virtuais (VMs - Virtual Machine). Esse processo é gerenciado por um componente denominado \textit{hipervisor}, que garante que cada máquina receba os recursos alocados sem interfirir no funcionamento de outras VMs. Cada ambiente virtual pode executar um sistema operacional próprio como se fossem máquinas separadas.

A Amazon Web Services (AWS), por exemplo, utiliza-se de infraestruturas virtualizadas para fornecer serviços a startups, grandes corporações e até órgãos governamentais, que desejam se utilizar dos benefícios da computação na nuvem de forma barata e escalável.

\section{Infraestrutura Escalável: A relação e as diferenças entre Clusters e Data Centers.}
Um Data Center é um local em que ficam guardados diversos equipamentos de TI, incluindo servidores. Os Clusters são agrupamentos de servidores conectados que trabalham juntos como se fossem um só.

Essa organização permite criar uma infraestrutura confiável e escalável, permitindo que os usuários aumentem ou diminuam a quantidade de recursos e usuários de acordo com suas necessidades

\section{Virtualização de Redes: O impacto e a arquitetura das abordagens SDN (Redes Definidas por Software).}

A virtualização de rede é um processo que combina todos os recursos físicos de rede para centralizar tarefas administrativas. Os administradores podem ajustar e controlar esses elementos virtualmente sem tocar nos seus componentes, o que simplifica muito o gerenciamento da rede.

Uma Rede definida por software(SDN – Software-Defined Networking) centraliza o roteamento de tráfego assumindo o gerenciamento de roteamento do roteamento de dados no ambiente físico. Por exemplo, você pode programar seu sistema para priorizar o tráfego de chamadas de vídeo sobre o tráfego de aplicações para garantir uma qualidade de chamada consistente em todas as reuniões online.

% Virtualização de funções de rede
% A tecnologia de virtualização de funções de rede combina as funções de dispositivos de rede, como firewalls, balanceadores de carga e analisadores de tráfego, que trabalham juntos para melhorar a performance da rede.

\section{Otimização e Desempenho: Tecnologias avançadas de aceleração de rede, como SR-IOV.} 

Em tipos de máquina virtual (VM) suportados, a Rede Acelerada permite Single Root I/O Virtualization (SR-IOV), ou virtualização de entrada e saída (E/S) de raiz única. Técnica que melhora significativamente o desempenho da rede. A rede acelerada permite que o tráfego de rede ignore a camada de comutador de software das pilhas de virtualização do Hyper-V, assim o tráfego de rede flui diretamente entre a VF e o hardware, fazendo com que a sobrecarga de E/S na camada de emulação de software diminua até que o desempenho da rede se torne quase o mesmo que os ambientes físicos locais.

\end{document}